\documentclass{article}
\usepackage{amssymb}
\usepackage{amsthm}
\usepackage{enumitem}
\usepackage{comment}
\usepackage{mathtools}
\usepackage{mathrsfs}
\theoremstyle{definition}
\newtheorem{definition}{Definition}[]
\theoremstyle{theorem}
\newtheorem{theorem}{Theorem}[]
\newtheorem{proposition}{Proposition}[]
\newtheorem{lemma}{Lemma}[]
\newtheorem{assumption}{Assumption}
\newtheorem{corollary}{Corollary}[]
\usepackage[utf8]{inputenc}
\usepackage{natbib}
\usepackage{parskip}
\usepackage{rsfso}
\usepackage{setspace}
\usepackage{subcaption}
\usepackage{xcolor,soul}

\providecommand{\myfloor}[1]{\left \lfloor #1 \right \rfloor }
\setlength{\tabcolsep}{10pt} % Default value: 6pt
\renewcommand{\arraystretch}{1.5} % Default value: 1


\usepackage{geometry}

 \geometry{
 a4paper,
 %total={170mm,257mm},
 left=35mm,
 right=35mm,
 top=30mm,
 }

\title{Optimal enfranchisement}
\author{Ulysse Lojkine\footnote{ulysse.lojkine[at]sciencespo.fr Sciences Po}}
\date{May 2026}

\onehalfspacing
\begin{document}

\maketitle

\begin{abstract}
 We study the optimal voting franchise for a decision which affects different persons with different intensities.
 \end{abstract}

A planner cares about a population $\mathcal{N} = \{ 1, ..., N \}$ where each agent is affected by a social binary decision (0 or 1). From the planner's point of view, agent $i$'s opinion is a Bernoulli random variable $X_i$.

To make a decision, the planner picks a franchise, i.e. a subset $C \subseteq \mathcal{N}$ of franchise members, or insiders. The other agents are nonmembers, or outsiders. The outcome under franchise $C$, denoted $Y_C$, is decided by a majority vote among members of $C$:
\begin{equation*}
Y_C = 1_{S_C > |C|/2} \text{ where }
S_C = \sum_{i \in C} X_i
\end{equation*}
Let us denote $p(i,C)$ the success probability of agent $i$ under franchise $C$:
\begin{equation*}
    p(i,C) = P[X_i = Y_C]
\end{equation*}
Because this strict majority rule is unbiased only if the number of voters is odd, we will assume throughout that $N$ is odd and consider only franchises of odd size.

Each agent's utility is denoted $U(i,C)$ and depends on her success. Aggregate utility under franchise $C$ is $U_C = \sum_{i=1}^N U(i,C)$, and the planner seeks to maximise its expected value $\mathbb{E}[U_C]$. Franchise $C'$ will be said to be superior to franchise $C$ if expected aggregate welfare is weakly higher under franchise $C'$ than under $C$, and franchise $C$ will be said to be optimal if it is superior to all franchises.

To sum up, given a population $\mathcal{N}$ and a joint probability distribution for $X_1, ..., X_N$ and $U(1,.), ..., U(N,.)$, the planner's problem is to find the corresponding optimal franchise or franchises. Since the set of possible franchises, i.e. subsets of $\mathcal{N}$ of odd size, is finite, such an optimal franchise always exists.

To make the problem tractable, let us make the following natural assumption:
\begin{assumption}[Selfishness]
\label{assumption:selfishness}
    Conditional on her own success, an agent is indifferent to the success of others, or to anyone's inclusion in the franchise.
\end{assumption}
In other words, the only way through which a franchise change affects an agent's utility is by determining whether or not she is successful. In particular, $\mathbb{E}[U(i,C)|X_i = Y_C]$ and $\mathbb{E}[U(i,C)|X_i \neq Y_C]$ do not vary with $C$, so that we can define:
\begin{equation*}
    \Delta_i = \mathbb{E}[U(i,C)|X_i = Y_C] - \mathbb{E}[U(i,C)|X_i \neq Y_C]
\end{equation*}
We call $\Delta_i$ the agent's \textit{stake} in the decision, and further assume it to be nonnegative -- i.e., the planner expects each agent to be on average weakly better off when successful. Without loss of generality, we will assume from now on that agents are ranked from the highest to the lowest stake: $\Delta_1 \geq  ... \geq \Delta_N \geq 0$. Then an individual's expected utility can be rewritten as:
\begin{equation*}
    \mathbb{E}[U(i,C)] = \Delta_i P[X_i = Y_C] + \mathbb{E}[U(i,C)|X_i \neq Y_C]
\end{equation*}


Also without loss of generality, we assume away the constant term in that expression, $\mathbb{E}[U(i,C)|X_i \neq Y_C]$, which is unaffected by the planner's decision. This allows to rewrite the planner's objective function as:
\begin{equation}
\label{eq:expected-aggregate-utility}
    \mathbb{E}[U_C] = \sum_{i \in \mathcal{N}} \Delta_i p(i,C)
\end{equation}

In other words, because only the expected welfare of the individuals matters to the planner, he doesn't have to care about the probability distribution of an individual's preference intensity. All that matters is its expectation $\Delta_i$ and the success probabilities $p(i,C)$. In practice, a planner would compute $\Delta_i$ as an expectation based on observable information about the potential voters, e.g. their place of residence, gender, age, etc. But formally, in all that follows, the problem is equivalent to one where the planner would know with certainty each agent's utility conditional on success.

A natural conjecture is that an optimal franchise would include all agents with a stake above a certain threshold, i.e., it would be a high-stake or a strictly high-stake franchise in the following sense.

\begin{definition}
    $C \subset \mathcal{N}$ is a (strictly) high-stake franchise if $C = \mathcal{N}$ or, for any $i \in C$ and $j \in \mathcal{N} \setminus C$, $\Delta_i$ is (strictly) larger than $\Delta_j$.
\end{definition}


We found even such a simple conjecture difficult to answer in the general case. For that reason, we will work from now on under two further assumptions on the \textit{ex ante} distribution of opinions.

\begin{assumption}[Uncertainty on opinions]
    The family of random variables $(X_i)_{i \in \mathcal{N}}$ has full support, i.e. for each $x \in \{0,1\}^N$, $P[(X_1, ..., X_n) = x] > 0$.
\end{assumption}

One way to interpret this assumption is that even conditional on knowing the opinion of all agents but one, the planner would still have some degree of uncertainty about the opinion of the last agent.

To motivate the exploration of optimal enfranchising, it can be useful to first show that such a procedure is superior to a planner's dictatorship, where the planner would use the information at his disposal to directly pick an outcome lottery independently of the agents' preferences.

\begin{definition}[Planner's dictatorship]
    Under a planner's dictatorship, the outcome is a Bernoulli variable $Z_q$ of parameter $q$ which is independent from $(X_1, ..., X_N)$s. The welfare of agent $i$ under this rule is denoted $V(i,q) = \Delta_i 1_{X_i = Z_q}$ and expected aggregate welfare is denoted $\mathbb{E}[V(q)] = \sum_{i \in \mathcal{N}} \Delta_i P[X_i = Z_q]$.
\end{definition}


\paragraph{Conjecture 1} There exists a franchise that provides weakly higher expected aggregate utility than any planner's dictatorship.

\paragraph{Conjecture 2} An optimal franchise is a high-stake franchise.

\end{document}
